% Options for packages loaded elsewhere
\PassOptionsToPackage{unicode}{hyperref}
\PassOptionsToPackage{hyphens}{url}
\documentclass[
  a4paper,
]{ltjsarticle}
\usepackage{xcolor}
\usepackage[margin=25mm]{geometry}
\usepackage{amsmath,amssymb}
\setcounter{secnumdepth}{-\maxdimen} % remove section numbering
\usepackage{iftex}
\ifPDFTeX
  \usepackage[T1]{fontenc}
  \usepackage[utf8]{inputenc}
  \usepackage{textcomp} % provide euro and other symbols
\else % if luatex or xetex
  \usepackage{unicode-math} % this also loads fontspec
  \defaultfontfeatures{Scale=MatchLowercase}
  \defaultfontfeatures[\rmfamily]{Ligatures=TeX,Scale=1}
\fi
\usepackage{lmodern}
\ifPDFTeX\else
  % xetex/luatex font selection
\fi
% Use upquote if available, for straight quotes in verbatim environments
\IfFileExists{upquote.sty}{\usepackage{upquote}}{}
\IfFileExists{microtype.sty}{% use microtype if available
  \usepackage[]{microtype}
  \UseMicrotypeSet[protrusion]{basicmath} % disable protrusion for tt fonts
}{}
\makeatletter
\@ifundefined{KOMAClassName}{% if non-KOMA class
  \IfFileExists{parskip.sty}{%
    \usepackage{parskip}
  }{% else
    \setlength{\parindent}{0pt}
    \setlength{\parskip}{6pt plus 2pt minus 1pt}}
}{% if KOMA class
  \KOMAoptions{parskip=half}}
\makeatother
\usepackage{longtable,booktabs,array}
\newcounter{none} % for unnumbered tables
\usepackage{calc} % for calculating minipage widths
% Correct order of tables after \paragraph or \subparagraph
\usepackage{etoolbox}
\makeatletter
\patchcmd\longtable{\par}{\if@noskipsec\mbox{}\fi\par}{}{}
\makeatother
% Allow footnotes in longtable head/foot
\IfFileExists{footnotehyper.sty}{\usepackage{footnotehyper}}{\usepackage{footnote}}
\makesavenoteenv{longtable}
\setlength{\emergencystretch}{3em} % prevent overfull lines
\providecommand{\tightlist}{%
  \setlength{\itemsep}{0pt}\setlength{\parskip}{0pt}}
\usepackage{bookmark}
\IfFileExists{xurl.sty}{\usepackage{xurl}}{} % add URL line breaks if available
\urlstyle{same}
\hypersetup{
  pdftitle={自己無撞着インフレーション機構の段階分解と N=3..40 全域再現 --- 再現仕様書（総括）},
  pdfauthor={木原 範昭（WF System Co., Ltd.）　ORCID 0009-0004-6753-4020},
  hidelinks,
  pdfcreator={LaTeX via pandoc}}

\title{自己無撞着インフレーション機構の段階分解と N=3..40 全域再現 ---
再現仕様書（総括）}
\author{木原 範昭（WF System Co., Ltd.）　ORCID 0009-0004-6753-4020}
\date{2026年9月5日　Version DOI 10.5281/zenodo.22317636}

\begin{document}
\maketitle

\textbf{著者:} 木原 範昭\\
\textbf{ORCID:} 0009-0004-6753-4020\\
\textbf{日付:} 2026年9月5日\\
\textbf{Version DOI:} 10.5281/zenodo.22317636\\
\textbf{Concept DOI:} 10.5281/zenodo.22317635（本総括と第1〜3章は本
Concept DOI を共有する）\\
\textbf{位置づけ:}
再現仕様書（総括＋章論文3本の構成）。物理的解釈は含まない

\begin{center}\rule{0.5\linewidth}{0.5pt}\end{center}

\subsection{1. 目的}\label{ux76eeux7684}

本仕様書群の目的は、\textbf{インフレーション的な発展------測定にかからないほど小さい種
（休眠比 H⊥/H \textasciitilde{}
10\^{}-30）が、力学だけで何十桁も指数増幅して飽和に至る現象------が起こる
機構を、再現可能な形で確定すること}である。

2026-07-22 の正本論文（自発的分裂の開始と帰結の三分類、Version DOI
10.5281/zenodo.21486234 / Concept DOI 10.5281/zenodo.21486233）で N=40,
300, 1000 に
観測されたこの現象について、力学の構成要素を一因子ずつ切り分けた結果、必要条件は
\textbf{段2（生成子の振幅正規化）＋段3（生成子の cos
対称部の除去＝実直交回転）}の同時成立で
あることが確定し（削除対照の監査記録は第2章補遺A）、その構成（段1+2+3
力学）を 静的親から N=3..40
の全域でスイープして普遍性を固定した（第2章）。

本総括は目的とアップロード構成の一覧のみを与える。\textbf{主章は第2章}であり、数式・
プログラム・データの対応、検証ゲート、観察はすべて章論文に記載されている。

\subsection{2. 論文構造（本 Concept DOI
配下）}\label{ux8ad6ux6587ux69cbux9020ux672c-concept-doi-ux914dux4e0b}

{\def\LTcaptype{none} % do not increment counter
\begin{longtable}[]{@{}
  >{\raggedright\arraybackslash}p{(\linewidth - 6\tabcolsep) * \real{0.2500}}
  >{\raggedright\arraybackslash}p{(\linewidth - 6\tabcolsep) * \real{0.2500}}
  >{\raggedright\arraybackslash}p{(\linewidth - 6\tabcolsep) * \real{0.2500}}
  >{\raggedright\arraybackslash}p{(\linewidth - 6\tabcolsep) * \real{0.2500}}@{}}
\toprule\noalign{}
\begin{minipage}[b]{\linewidth}\raggedright
章
\end{minipage} & \begin{minipage}[b]{\linewidth}\raggedright
ファイル
\end{minipage} & \begin{minipage}[b]{\linewidth}\raggedright
内容
\end{minipage} & \begin{minipage}[b]{\linewidth}\raggedright
式番号
\end{minipage} \\
\midrule\noalign{}
\endhead
\bottomrule\noalign{}
\endlastfoot
総括（本書） & \texttt{paper\_overview/overview\_stage123\_sweep\_ja.md}
& 目的・構成一覧・同梱物台帳 & --- \\
第1章 & \texttt{paper\_ch1\_static\_parents/ch1\_static\_parents\_ja.md}
& 静的親データの生成（rng式・make\_parent・零閉塞種・Z0、7月正本との bit
一致ゲート） & 式1〜式9 \\
\textbf{第2章（主章）} &
\texttt{paper\_ch2\_sweep\_dynamics/ch2\_sweep\_stage123\_ja.md} &
段1+2+3
力学の定義（旧力学の数学・新旧回転写像の差異・分母設計仮説・休眠比の指標性・複素図との対応）と
N=3..40 スイープ。補遺A=削除対照監査表、補遺B=プログラム系譜 &
式10〜式25 \\
第3章 & \texttt{paper\_ch3\_complex\_plane/ch3\_complex\_plane\_ja.md} &
複素平面読出し図3種（step0・終了時・凝縮中心拡大）の実装と観察 &
式26〜式28 \\
\end{longtable}
}

\subsection{3. Zenodo
アップロード一覧}\label{zenodo-ux30a2ux30c3ux30d7ux30edux30fcux30c9ux4e00ux89a7}

\subsubsection{3.1 論文}\label{ux8ad6ux6587}

\begin{itemize}
\tightlist
\item
  総括・第1章・第2章・第3章の md（日本語正本）および
  tex/PDF（日英、Zenodo 公開時に生成）
\end{itemize}

\subsubsection{\texorpdfstring{3.2 本体パッケージ
\texttt{N3\_N40\_stage123\_sweep\_20260905/}（約476MB）}{3.2 本体パッケージ N3\_N40\_stage123\_sweep\_20260905/（約476MB）}}\label{ux672cux4f53ux30d1ux30c3ux30b1ux30fcux30b8-n3_n40_stage123_sweep_20260905ux7d04476mb}

{\def\LTcaptype{none} % do not increment counter
\begin{longtable}[]{@{}
  >{\raggedright\arraybackslash}p{(\linewidth - 2\tabcolsep) * \real{0.5000}}
  >{\raggedright\arraybackslash}p{(\linewidth - 2\tabcolsep) * \real{0.5000}}@{}}
\toprule\noalign{}
\begin{minipage}[b]{\linewidth}\raggedright
区分
\end{minipage} & \begin{minipage}[b]{\linewidth}\raggedright
内容
\end{minipage} \\
\midrule\noalign{}
\endhead
\bottomrule\noalign{}
\endlastfoot
プログラム（6本） &
\texttt{make\_static\_parents\_N3\_N40\_v1.py}（親生成）・\texttt{run\_N3\_N40\_stage123\_v1.py}（スイープ本体）・\texttt{check\_sweep\_inputs\_v1.py}（入力ゲート）・\texttt{analyze\_sweep\_summary\_v1.py}（集計）・\texttt{plot\_complex\_plane\_N3\_N40\_stage123\_v1.py}（図化）・\texttt{run\_n\_scaling\_lowrank\_v1.py}（7月正本エンジンの
bit 同一コピー） \\
一括再現 & \texttt{run\_all.sh}（親生成→スイープ→ゲート→集計→図化） \\
初期データ & \texttt{parents/} 静的親 npz 38個＋台帳
\texttt{parents\_summary.csv}（約636KB） \\
結果データ & \texttt{results/} 状態 npz
228個（全501step×全波保存）・timeseries CSV（114,228行）・summary
CSV（228行）・集計 JSON・RUN\_METADATA \\
図（4枚） & H⊥/H
分母コントロール図（目標図）・複素平面3グリッド（step0/終了時/拡大） \\
台帳 &
\texttt{README.md}（進行記録・章登録）・\texttt{SHA256SUMS.txt}（全ファイルの
SHA256 正本） \\
\end{longtable}
}

\subsubsection{3.3
再現性担保のための関連フォルダ（補遺A・Bの実物）}\label{ux518dux73feux6027ux62c5ux4fddux306eux305fux3081ux306eux95a2ux9023ux30d5ux30a9ux30ebux30c0ux88dcux907aabux306eux5b9fux7269}

{\def\LTcaptype{none} % do not increment counter
\begin{longtable}[]{@{}
  >{\raggedright\arraybackslash}p{(\linewidth - 4\tabcolsep) * \real{0.3333}}
  >{\raggedright\arraybackslash}p{(\linewidth - 4\tabcolsep) * \real{0.3333}}
  >{\raggedright\arraybackslash}p{(\linewidth - 4\tabcolsep) * \real{0.3333}}@{}}
\toprule\noalign{}
\begin{minipage}[b]{\linewidth}\raggedright
フォルダ
\end{minipage} & \begin{minipage}[b]{\linewidth}\raggedright
内容
\end{minipage} & \begin{minipage}[b]{\linewidth}\raggedright
容量
\end{minipage} \\
\midrule\noalign{}
\endhead
\bottomrule\noalign{}
\endlastfoot
\texttt{ChatGPT\_denominator\_controls\_N40\_selfcontrol\_20260904/} &
段構成を確定した N=40 一因子実験の全実物（62ファイル）:
自己対照・静的親差替（段1基準）・段2・段3・段2削除・段2初期化のみ・σ時計の各プログラム／results／図／README／SHA256SUMS
& 約238MB \\
\texttt{自発的分裂予備実験\_v1\_N40対照実験系\_20260904/} & 7月正本 N=40
走行のゲート付き再現（fcurve 全行 bit 一致）・正本静的親
\texttt{parent\_static\_N40\_makeparent\_20260904.npz}・複素図・インフレーション図
& 約1.5MB \\
\end{longtable}
}

\subsubsection{3.4
引用で参照する既公開物（再アップロードしない）}\label{ux5f15ux7528ux3067ux53c2ux7167ux3059ux308bux65e2ux516cux958bux7269ux518dux30a2ux30c3ux30d7ux30edux30fcux30c9ux3057ux306aux3044}

\begin{itemize}
\tightlist
\item
  7月正本論文・プログラム・図:
  「N体関係波閉鎖系における状態の自発的分裂の開始と帰結の
  三分類」v1.0（2026-07-22）。\textbf{Version DOI
  10.5281/zenodo.21486234 / Concept DOI
  10.5281/zenodo.21486233}。レコードに再現プログラム
  zip（\texttt{run\_spontaneous\_splitting\_largeN\_v1.py}・
  エンジン等）と図が同梱済み。
\end{itemize}

\subsection{4. 検証の要 ---
全体を貫くゲート連鎖}\label{ux691cux8a3cux306eux8981-ux5168ux4f53ux3092ux8cabux304fux30b2ux30fcux30c8ux9023ux9396}

本仕様書群の再現性は、次の bit レベルの連鎖で7月正本に接続されている
（各ゲートの定義・実測は該当章）:

\begin{enumerate}
\def\labelenumi{\arabic{enumi}.}
\tightlist
\item
  \textbf{7月正本 ⇔ 再現走行}:
  \texttt{自発的分裂予備実験\_v1\_N40対照実験系\_20260904}
  の無変更再走行が 7月 fcurve と全 3512 行 bit 一致（GATE1/GATE2）
\item
  \textbf{再現走行 ⇔ 静的親}: 正本静的親の v・g・Z0 が7月走行の初期値と
  bit 一致（第1章 G1）。 第1章で生成した N=40 親は正本静的親と bit
  一致（npz の SHA256 も同一）
\item
  \textbf{静的親 ⇔ スイープ}: 全228走行の保存済み Z{[}0{]} が各 N
  の静的親 Z0 と bit 一致 （第2章 G1、checked 228 / MISMATCH 0）
\end{enumerate}

全ファイルの SHA256 は各パッケージ同梱の \texttt{SHA256SUMS.txt}
を正本とする。

\subsection{5.
実行環境（共通）}\label{ux5b9fux884cux74b0ux5883ux5171ux901a}

\begin{itemize}
\tightlist
\item
  Python 3.9.6（venv）・numpy 2.0.2（BLAS/LAPACK: macOS
  Accelerate）・matplotlib
\item
  macOS 26.3.1（arm64, Darwin 25.x）・乱数 numpy
  PCG64（\texttt{default\_rng}）
\end{itemize}

\begin{center}\rule{0.5\linewidth}{0.5pt}\end{center}

（総括おわり。詳細はすべて第1〜3章および同梱 README・SHA256SUMS を参照）

\end{document}
